\documentclass[12pt]{article}

% ── Packages ──────────────────────────────────────────────────────────────────
\usepackage[margin=1in]{geometry}
\usepackage{times}
\usepackage[T1]{fontenc}
\usepackage[utf8]{inputenc}
\usepackage{microtype}
\usepackage{setspace}
\usepackage{titlesec}
\usepackage{booktabs}
\usepackage{array}
\usepackage{graphicx}
\usepackage{xcolor}
\usepackage{hyperref}
\usepackage{amsmath}
\usepackage{parskip}
\usepackage{caption}
\usepackage{subcaption}
\usepackage{float}
\usepackage{enumitem}
\usepackage{fancyhdr}
\usepackage{natbib}

% ── Page style ────────────────────────────────────────────────────────────────
\pagestyle{fancy}
\fancyhf{}
\rhead{}
\lhead{\small Liakh}
\cfoot{\thepage}
\renewcommand{\headrulewidth}{0.4pt}

% ── Section formatting ────────────────────────────────────────────────────────
\titleformat{\section}{\large\bfseries}{\thesection}{1em}{}
\titleformat{\subsection}{\normalsize\bfseries}{\thesubsection}{1em}{}
\titleformat{\subsubsection}{\normalsize\itshape}{\thesubsubsection}{1em}{}

% ── Hyperref setup ────────────────────────────────────────────────────────────
\hypersetup{
  colorlinks=true,
  linkcolor=blue!60!black,
  citecolor=blue!60!black,
  urlcolor=blue!60!black
}

% ── Spacing ───────────────────────────────────────────────────────────────────
\setstretch{1.15}
\setlength{\parindent}{0pt}
\setlength{\parskip}{8pt}

% ══════════════════════════════════════════════════════════════════════════════
\begin{document}

% ── Title page ────────────────────────────────────────────────────────────────
\begin{titlepage}
  \centering
  \vspace*{4cm}

  {\LARGE\bfseries When Detection Fails, Fixing Doesn't Matter:\\[0.4em]
  Composing LLMs for Automated Security Patching\par}

  \vspace{0.5cm}
  {\large\itshape Can LLMs replace static analysis for vulnerability detection and repair?\\[0.2em]
  An empirical evaluation across 7 models, 49 pipeline combinations, and 22 CWE categories\par}

  \vspace{1.2cm}
  {\normalsize
    Timothy Liakh\\[0.15em]
    Rensselaer Polytechnic Institute\\[0.15em]
    Advisor: Konstantin Kuzmin\\[0.15em]
    Summer 2026
  \par}

  \vspace{0.5cm}
  {\normalsize Repository: \href{https://github.com/Vitisadi/SecurePatch}{\texttt{github.com/Vitisadi/SecurePatch}}\par}
\end{titlepage}

% ── Abstract ──────────────────────────────────────────────────────────────────
\newpage
\begin{center}
  \textbf{Abstract}
\end{center}

Static analysis tools have long been the standard for catching security
vulnerabilities before code ships. They are fast and consistent, but they
struggle with context: they flag patterns, not intent, and they cannot generate
a fix. Large language models offer a different trade-off. They reason about code
semantically, can explain why something is wrong, and can propose a patch in the
same step. Whether they are actually better than static analysis at finding real
vulnerabilities, and in which CWE categories they succeed or fail, is an open
question this project tries to answer empirically.

The study evaluates seven LLMs as vulnerability detectors across 56 real
cases spanning 22 CWE categories, drawn from CWEval (44 cases), hand-seeded
examples (6), and published literature (6). Detection and repair are treated as
separate, composable steps. A detector scans code and produces findings; a fixer
receives those findings and generates a patch. Running all combinations of the
seven detectors and seven fixers produces a 7$\times$7 matrix of 49 pipelines,
each evaluated on fix rate, regression rate, F1, and cost.

Solo models peak at 75\,\% F1 (GPT-5.5). A ground-truth experiment that hands
each fixer a perfect detection with zero false positives shows Opus reaching
76\,\% fix rate in isolation. The gap between that ceiling and real pipeline
performance traces almost entirely to false positives in the detection step.
Across the heatmap, the GPT-5.5 detector row produces the best F1 in most
fixer columns (up to 75\,\%), while Ollama lags at 28--50\,\%.

Cost varies by three orders of magnitude. Haiku detection paired with
GPT-4.1-mini fixing reaches 68\,\% F1 for \$0.11 across all 56 cases. The
GPT-5.5 self-pipeline costs \$7.10 for 75\,\% F1. Ensemble detection using a
majority vote across frontier models pushes F1 to 92\,\%, at the cost of
running multiple detectors per case.

The central finding is that detector false-positive rate determines the ceiling
of any pipeline. Upgrading the fixer without addressing detection quality
produces diminishing returns.

% ── Table of contents ─────────────────────────────────────────────────────────
\newpage
\tableofcontents

% ══════════════════════════════════════════════════════════════════════════════
\newpage
\section{Introduction}

Static analysis security testing has been the default tool for catching
vulnerabilities in production code for over two decades. Tools like Semgrep,
CodeQL, and Bandit operate on fixed rule sets: they are fast, deterministic,
and produce no false negatives for the patterns they cover. The problem is that
coverage is the ceiling. A rule that does not exist catches nothing, and writing
rules requires security expertise most development teams do not have. More
importantly, none of these tools fix what they find. A developer still has to
read the alert, understand the root cause, and write the patch.

Large language models change the equation on both sides. They are not bound by
explicit rules; they reason about code semantically and can generalize across
vulnerability classes without being programmed for each one. They can also
generate a patch in the same pass. But whether that reasoning is actually
reliable enough to replace or augment static analysis, and which models perform
well at detection versus repair, has not been studied systematically across a
controlled benchmark.

This project measures exactly that. Seven LLMs are evaluated as detectors and
as fixers, independently, across 56 real vulnerabilities spanning 22 CWE
categories. Detection and repair are treated as composable stages: a detector
produces findings, and a fixer receives those findings and generates a patch.
Every combination of the seven detectors and seven fixers is run, producing 49
distinct pipelines. Each pipeline is judged by a deterministic oracle on fix
rate, regression rate, F1, and cost.

The central question is:

\medskip
\noindent\textbf{RQ1.}\quad\textit{Can LLMs match or exceed static analysis
for vulnerability detection, and does composing a strong detector with a strong
fixer produce better security outcomes than either model working alone?}

\smallskip
\noindent\textbf{RQ2.}\quad\textit{Does running multiple detectors and taking
a majority vote meaningfully improve detection quality, and is the added cost
justified by the gain in F1?}
\medskip

\subsection{Contributions}

\begin{enumerate}[leftmargin=*, label=\arabic*.]

  \item \textbf{A controlled vulnerability benchmark.} 56 cases drawn from
  CWEval (44), hand-seeded examples (6), and published literature (6), covering
  22 CWE categories. Each case has a ground-truth vulnerable file, a known fix,
  and a deterministic oracle that distinguishes a correct patch from a
  regression or a no-op.

  \item \textbf{A reproducible evaluation harness.} The
  \texttt{CachedFindingDetector} replays pre-computed detection outputs
  identically across all fixer columns, ensuring that differences in pipeline
  F1 reflect fixer behavior rather than detection variance. All results are
  stored as versioned JSONL files and can be regenerated from the repository.

  \item \textbf{A full 7$\times$7 pipeline heatmap.} Every combination of
  seven detectors and seven fixers is evaluated, producing 49 F1 scores on the
  same 56-case benchmark. This is the primary empirical contribution of the
  project.

  \item \textbf{A ground-truth fixer experiment.} Each fixer is given a
  perfect detection with zero false positives and 100\,\% recall, isolating
  pure repair capability from detection noise. The gap between this ceiling and
  real pipeline F1 quantifies how much detection quality costs each fixer.

  \item \textbf{An ensemble detection study.} Majority-vote combinations of
  multiple detectors are evaluated on F1, recall, and false-positive count.
  The best configuration, a frontier majority vote across four models, reaches
  92\,\% F1.

\end{enumerate}

\subsection{Roadmap}

The paper begins by grounding this work relative to existing research on
LLM-based vulnerability detection and automated program repair. It then
describes how the benchmark was constructed and how the evaluation pipeline
works. The results follow in order of increasing complexity: solo model
performance first, then the full composition matrix, then cost and ensemble
analysis. The paper closes with a discussion of what the findings mean in
practice, the limits of the current study, and directions for follow-on work.

% ══════════════════════════════════════════════════════════════════════════════
\section{Background and Related Work}
\label{sec:background}

\subsection{LLM-Based Vulnerability Detection}

Detecting security vulnerabilities in source code has traditionally been the
domain of static analysis. Tools like Semgrep and CodeQL match code against
hand-written rules, offering speed and determinism at the cost of coverage:
they catch only what their rules describe. Machine learning approaches tried to
close that gap, with models trained on historical vulnerability data to predict
whether a function or line was dangerous~\cite{fu2022linevul, thapa2022transformer, li2023sast}.
These models improved generalization but introduced a new problem: they produce
false positives at rates that make triage expensive in practice.

Large language models enter this space differently. Rather than classifying
against learned patterns, they read code and reason about what it does. Pearce
et al.\ showed that LLMs could identify security weaknesses in code generated
by GitHub Copilot~\cite{pearce2022copilot}, and later demonstrated that
models could detect and explain vulnerabilities in hand-crafted scenarios
without any task-specific fine-tuning~\cite{pearce2023repair}. Meta's
CyberSecEval benchmark evaluated frontier models on a range of security tasks
and found substantial variance across CWE categories~\cite{cyberSecEval}.
What these studies share is a focus on a single model acting alone. This
project asks a different question: what happens when you pick one model to
detect and a different model to fix?

\subsection{Automated Program Repair}

Automated program repair has a long history independent of security. The
core idea is that given a failing test and a buggy program, a system should
be able to produce a patch without human intervention. Early approaches used
genetic algorithms and constraint-based synthesis. LLMs changed what was
possible. Xia et al.\ found that pre-trained code models could generate
correct patches at rates that surpassed prior symbolic methods on standard
benchmarks~\cite{xia2023apr}, and Jiang et al.\ showed that model scale was
a reliable predictor of repair success~\cite{jiang2023apr}. Sobania et al.\
evaluated ChatGPT specifically on the QuixBugs benchmark and found it
competitive with dedicated APR tools on simple bugs~\cite{sobania2023chatgpt}.

The security repair problem is harder than general APR. A patch that makes a
test pass is not necessarily a secure patch: it may suppress the symptom
without removing the vulnerability, or introduce a regression elsewhere.
Pearce et al.\ examined this directly, finding that LLMs often generated
plausible-looking but insecure fixes when prompted without explicit security
guidance~\cite{pearce2023repair}. This project addresses that by using a
deterministic oracle that checks both fix success and regression, and by
separating the detection signal from the repair step so the fixer receives
an explicit vulnerability description rather than just a failing test.

\subsection{Existing Benchmarks}

SWE-bench~\cite{swebench} is the closest large-scale benchmark to this
project. It draws real GitHub issues from Python repositories and validates
patches against the project's own test suite. The focus there is model
capability on general software tasks. SecurityEval~\cite{securityeval}
targets security specifically, using 130 code generation prompts mapped to
CWE categories to test whether models produce vulnerable output. CyberSecEval
takes a broader view, covering prompt injection, insecure code generation,
and vulnerability identification across multiple languages~\cite{cyberSecEval}.

None of these benchmarks treat detection and repair as a composed pipeline.
They either test one capability at a time or evaluate a single model end to
end. The 7$\times$7 design here is, to our knowledge, the first systematic
study of cross-model composition for security vulnerability repair.

\subsection{Where This Project Fits}

This project is closest in spirit to SWE-bench but more controlled in design.
Every case has a known-vulnerable commit, a verified fix, and a deterministic
oracle that does not rely on an LLM judge. The central
contribution is not a new model or a new benchmark in isolation, but a
methodology for measuring what each stage of a detect-then-fix pipeline
actually contributes to the final outcome.

% ══════════════════════════════════════════════════════════════════════════════
\section{Benchmark Design}
\label{sec:benchmark}

\subsection{Case Selection and Dataset Composition}

The benchmark contains 56 cases drawn from three sources. The majority (44)
come from CWEval, an open collection of vulnerability scenarios designed for
evaluating LLM security capabilities. Six additional cases were hand-seeded:
written from scratch to cover CWE categories underrepresented in CWEval or
to test boundary conditions that naturally occurring examples rarely produce.
The remaining six come from published vulnerability disclosures and literature
examples where a real fix commit is publicly available.

Every case provides the same inputs to the pipeline: a source file containing
a vulnerability, the CWE category it belongs to, and a reference fix. Cases
were selected to have self-contained vulnerabilities, meaning the flaw and its
fix are localized to a single file. This keeps the evaluation clean: a fixer
that modifies the right file either produces a correct patch or it does not,
and the oracle can verify this without reasoning about cross-file dependencies.

\subsection{CWE Coverage}

The 56 cases span 22 distinct CWE categories. Table~\ref{tab:cwes} lists each
CWE present in the benchmark along with its case count. The distribution is
intentionally varied: no single category dominates, which prevents a model
from scoring well by specializing in one vulnerability class. The most
represented categories reflect the types of flaws most commonly found in
real code, including injection vulnerabilities, improper input validation,
and cryptographic weaknesses.

\begin{table}[H]
  \centering
  \caption{CWE categories covered in the benchmark (56 cases total, 22 CWEs).}
  \label{tab:cwes}
  \small
  \begin{tabular}{llc}
    \toprule
    CWE ID & Description & Cases \\
    \midrule
    CWE-78  & OS Command Injection                        & 4 \\
    CWE-79  & Cross-site Scripting (XSS)                  & 3 \\
    CWE-89  & SQL Injection                                & 4 \\
    CWE-22  & Path Traversal                               & 3 \\
    CWE-94  & Code Injection                               & 2 \\
    CWE-327 & Use of Broken/Risky Cryptographic Algorithm  & 3 \\
    CWE-330 & Use of Insufficiently Random Values          & 2 \\
    CWE-502 & Deserialization of Untrusted Data            & 2 \\
    CWE-611 & XML External Entity (XXE)                    & 2 \\
    CWE-798 & Use of Hard-coded Credentials                & 2 \\
    CWE-20  & Improper Input Validation                    & 4 \\
    CWE-400 & Uncontrolled Resource Consumption            & 2 \\
    CWE-434 & Unrestricted File Upload                     & 2 \\
    CWE-601 & Open Redirect                                & 2 \\
    CWE-918 & Server-Side Request Forgery (SSRF)           & 2 \\
    CWE-295 & Improper Certificate Validation              & 2 \\
    CWE-306 & Missing Authentication                       & 2 \\
    CWE-352 & Cross-Site Request Forgery (CSRF)            & 2 \\
    CWE-476 & NULL Pointer Dereference                     & 2 \\
    CWE-190 & Integer Overflow                             & 2 \\
    CWE-312 & Cleartext Storage of Sensitive Information   & 2 \\
    CWE-416 & Use After Free                               & 2 \\
    \midrule
    \multicolumn{2}{l}{\textit{Total}} & 56 \\
    \bottomrule
  \end{tabular}
\end{table}

\subsection{Evaluation Oracle}

Each case is judged by a deterministic oracle rather than an LLM reviewer.
The oracle runs in a Docker container to ensure consistent execution across
all pipeline combinations. Given the original vulnerable file and the fixer's
output, it checks three conditions in order: whether the patched file compiles
or parses correctly, whether the targeted vulnerability is no longer present
according to a rule matched to the case's CWE, and whether the fix introduced
any regression by running a suite of behavioral checks against the original
file's expected outputs.

A case is counted as a fix only if all three conditions pass. A patch that
removes the vulnerability but breaks existing behavior is recorded as a
regression, not a fix. A patch that makes no change, or changes code
unrelated to the vulnerability, is recorded as a no-op. This three-way
verdict is what allows fix rate and regression rate to be reported
independently rather than collapsed into a single success metric.

% ══════════════════════════════════════════════════════════════════════════════
\section{Pipeline Architecture and Methodology}
\label{sec:methodology}

\subsection{Detect $\rightarrow$ Enrich $\rightarrow$ Fix Pipeline}

% Figure 2 placeholder
\begin{figure}[H]
  \centering
  \fbox{\parbox{0.7\textwidth}{\centering\vspace{2cm}
    \textit{Figure 2: Pipeline diagram — Detector $\rightarrow$ cached findings
    $\rightarrow$ Fixer $\rightarrow$ oracle verdict}
  \vspace{2cm}}}
  \caption{End-to-end pipeline architecture.}
  \label{fig:pipeline}
\end{figure}

\subsection{CachedFindingDetector}

\subsection{Models Under Test}

% Table 1 placeholder
\begin{table}[H]
  \centering
  \caption{Models used as detectors and fixers.}
  \label{tab:models}
  \begin{tabular}{llll}
    \toprule
    Model & Provider & Tier & Role(s) \\
    \midrule
    Claude Sonnet 4.6   & Anthropic & Frontier & Detector, Fixer \\
    Claude Opus 4.8     & Anthropic & Frontier & Detector, Fixer \\
    Claude Haiku 4.5    & Anthropic & Budget   & Detector, Fixer \\
    GPT-5.5             & OpenAI    & Frontier & Detector, Fixer \\
    GPT-4.1-mini        & OpenAI    & Budget   & Detector, Fixer \\
    Gemini 2.5 Flash    & Google    & Frontier & Detector, Fixer \\
    Ollama 7b (local)   & Open      & Local    & Detector, Fixer \\
    \bottomrule
  \end{tabular}
\end{table}

\subsection{Metrics}

The primary metric is F1, computed as:
\[
  F1 = \frac{2 \cdot \mathrm{TP}}{2 \cdot \mathrm{TP} + \mathrm{FP} + \mathrm{FN}}
\]
where TP is the number of cases fixed without regression, FP is the number of
detector false positives across all 56 cases, and FN is $56 - \mathrm{TP}$.
Fix rate, regression rate, cost (USD per 56-case run), and fixes per dollar are
also reported.

\subsection{Experimental Controls and Methodology Corrections}

% ══════════════════════════════════════════════════════════════════════════════
\section{Results}
\label{sec:results}

\subsection{Solo Pipeline Performance}

% Table 2 placeholder
\begin{table}[H]
  \centering
  \caption{Solo pipeline results (detector = fixer, diagonal of heatmap).}
  \label{tab:solo}
  \begin{tabular}{lrrrrr}
    \toprule
    Model & Fix Rate & F1 & Regressions & Cost (\$) & Fixes/\$ \\
    \midrule
    Opus 4.8       & 66\% & 70\% & 10 & 0.73  & 51   \\
    GPT-5.5        & 64\% & 75\% & 11 & 1.84  & 20   \\
    Sonnet 4.6     & 62\% & 67\% & 18 & 0.33  & 106  \\
    GPT-4.1-mini   & 54\% & 63\% & 13 & 0.03  & 1071 \\
    Gemini 2.5     & 50\% & 56\% & 23 & 0.05  & 596  \\
    Haiku 4.5      & 59\% & 62\% & 13 & 0.11  & 306  \\
    Ollama 7b      & 23\% & 28\% & 35 & 0.00  & free \\
    \bottomrule
  \end{tabular}
\end{table}

\subsection{Full 7$\times$7 Pipeline Heatmap}

% Figure 3 placeholder
\begin{figure}[H]
  \centering
  \fbox{\parbox{0.85\textwidth}{\centering\vspace{3cm}
    \textit{Figure 3: 7$\times$7 F1 heatmap (detector rows $\times$ fixer columns)}
  \vspace{3cm}}}
  \caption{F1 scores across all 49 detector-fixer combinations. Color scale:
  green = high F1, red = low.}
  \label{fig:heatmap}
\end{figure}

\subsection{Detector Quality as the Dominant Variable}

% Figure 4 placeholder
\begin{figure}[H]
  \centering
  \fbox{\parbox{0.7\textwidth}{\centering\vspace{2cm}
    \textit{Figure 4: Scatter — detector FP count vs.\ row-best F1}
  \vspace{2cm}}}
  \caption{Each point is one detector. Fewer false positives correlate with
  higher achievable F1 regardless of which fixer is paired.}
  \label{fig:fp-scatter}
\end{figure}

\subsection{Ground-Truth Fixer Experiment}

% Table 3 placeholder
\begin{table}[H]
  \centering
  \caption{Fixer performance given perfect detection (56/56 recall, 0 FPs).}
  \label{tab:gt}
  \begin{tabular}{lrr}
    \toprule
    Fixer & Fixed & Fix Rate \\
    \midrule
    Opus 4.8      & 43/56 & 76\% \\
    GPT-5.5       & 40/56 & 71\% \\
    Haiku 4.5     & 38/56 & 67\% \\
    Sonnet 4.6    & 36/56 & 64\% \\
    GPT-4.1-mini  & 35/56 & 62\% \\
    Gemini 2.5    & 28/56 & 50\% \\
    Ollama 7b     & \multicolumn{2}{c}{all errors} \\
    \bottomrule
  \end{tabular}
\end{table}

\subsection{Cost-Efficiency Analysis}

% Figure 5 placeholder
\begin{figure}[H]
  \centering
  \fbox{\parbox{0.7\textwidth}{\centering\vspace{2cm}
    \textit{Figure 5: F1 vs.\ combined pipeline cost — each of 49 combos plotted,
    Pareto frontier highlighted}
  \vspace{2cm}}}
  \caption{Cost-efficiency of all 49 pipeline combinations. Points on the Pareto
  frontier deliver the highest F1 at a given budget.}
  \label{fig:pareto}
\end{figure}

\subsection{Ensemble Detection}

% Table 4 placeholder
\begin{table}[H]
  \centering
  \caption{Ensemble detection configurations.}
  \label{tab:ensemble}
  \begin{tabular}{llrrr}
    \toprule
    Strategy & Models & F1 & Recall & FPs \\
    \midrule
    Frontier $\geq$2 & Sonnet + Opus + GPT-5.5 + Gemini & 92\% & 96\% & 8  \\
    Claude $\geq$2   & Sonnet + Opus + Haiku            & 90\% & 98\% & 11 \\
    OpenAI $\geq$1   & GPT-5.5 + GPT-mini               & 87\% & 93\% & 12 \\
    Budget $\geq$1   & GPT-mini + Haiku                 & 81\% & 96\% & 23 \\
    Union all 7      & All models                       & 67\% & 100\% & -- \\
    \bottomrule
  \end{tabular}
\end{table}

% ══════════════════════════════════════════════════════════════════════════════
\section{Discussion}
\label{sec:discussion}

\subsection{Why the Detector Matters More Than the Fixer}

\subsection{The Haiku Fixer Failure}

\subsection{Practical Pairing Recommendations}

\subsection{What Ensemble Detection Buys You}

% ══════════════════════════════════════════════════════════════════════════════
\section{Threats to Validity}
\label{sec:threats}

\subsection{Construct Validity}

\subsection{Internal Validity}

\subsection{External Validity}

\subsection{Reproducibility}

% ══════════════════════════════════════════════════════════════════════════════
\section{Future Work}
\label{sec:future}

% ══════════════════════════════════════════════════════════════════════════════
\section{Conclusion}

% ── References ────────────────────────────────────────────────────────────────
\newpage
\bibliographystyle{plainnat}
\bibliography{references}

% ── Appendices ────────────────────────────────────────────────────────────────
\newpage
\appendix

\section{Full Heatmap Data}

\begin{table}[H]
  \centering
  \caption{F1 (\%) for all 49 detector-fixer combinations (v2 methodology).}
  \label{tab:full-heatmap}
  \small
  \begin{tabular}{lrrrrrrr}
    \toprule
    Det $\backslash$ Fix & Sonnet & Opus & Haiku & GPT-5.5 & Mini & Gemini & Ollama \\
    \midrule
    Sonnet   & 67 & 69 & 61 & 70 & 67 & 58 & 35 \\
    Opus     & 68 & 70 & 65 & 70 & 65 & 58 & 40 \\
    Haiku    & 70 & 71 & 62 & 69 & 68 & 60 & 42 \\
    GPT-5.5  & 72 & 75 & 74 & 75 & 74 & 61 & 42 \\
    Mini     & 62 & 67 & 69 & 69 & 63 & 49 & 40 \\
    Gemini   & 65 & 68 & 63 & 69 & 63 & 56 & 42 \\
    Ollama   & 45 & 50 & 59 & 50 & 46 & 40 & 28 \\
    \bottomrule
  \end{tabular}
\end{table}

\section{Per-Case Breakdown}

% Table to be populated from harness/results/*.jsonl

\end{document}
